\paragraph{C.3ba (Refuted): integral finiteness in positive characteristic.}
Let $k$ be any algebraically closed field of characteristic $p>0$.
For every pair of integers $N,r$ with $2\le r\le N-1$, there is
an explicitly specified excellent normal strictly henselian local
$k$-algebra $R'$ of dimension $N$ and residue field $k$, with regular
puncture $X'$, for which the ordinary quotient admits a surjection
\[
 B^r(X')=
 \mathrm{CH}^r(X')/\mathrm{CH}^r(X')_{\mathrm{alg}}
                         \longrightarrow (k,+).
\]
Thus integral finiteness is false throughout this range, including
over $\overline{\mathbb F}_p$. The exhibited quotient has exponent
$p$ and proves no nonzero rational quotient.

\begin{proof}
\emph{Step 1: specified projective and strict-local data.}
Use C.3as with projective dimension $d=r$. It supplies an explicitly
constructed smooth connected projective $r$-fold $V$, an ample line
bundle $L$, and a nowhere-zero section
\[
 \eta\in H^0(V,\Omega^1_{V/k}\otimes L^{-1}).
\]
Write
\[
 A=\bigoplus_{n\ge0}H^0(V,L^n),\quad
 R=(A_{A_+})^{\mathrm{sh}},\quad
 Z_0=\operatorname{Tot}(L^{-1}),\quad
 Z=Z_0\times_{\operatorname{Spec}A}\operatorname{Spec}R.
\]
Its reduced exceptional divisor is $V$. C.3as constructs actual
closed curves $D_b\subset Z$, one for each $b\in k$, meeting $V$
only at the fixed point $v$. In the fixed etale coordinates
$s_1,\ldots,s_{r-1},z$ at $v$ and the fibre coordinate $t$, their
ideals near $v$ are
\[
                  (s_1,\ldots,s_{r-1},z-bt).
\]
The normalized form is $\eta=dt^{-1}\otimes dz$, and the integral
conormal invariant of C.3as satisfies
$\rho([D_b]-[D_0])=b$. These curves are the individually selected
henselian branches proved there; no formal cycles are used.

Put $m=N-r-1$. If $m=0$, this is exactly C.3as. Suppose henceforth
that $m\ge1$. With coordinates $[z_0:\cdots:z_m]$ set
\[
 V'=\mathbb P^m\times V,\qquad
 L'=\mathcal O_{\mathbb P^m}(1)\boxtimes L,
\]
and define
\[
 \begin{aligned}
 A'&=\bigoplus_{n\ge0}H^0(V',(L')^n),&
 R'&=(A'_{A'_+})^{\mathrm{sh}},\\
 Z'_0&=\operatorname{Tot}((L')^{-1}),&
 Z'&=Z'_0\times_{\operatorname{Spec}A'}\operatorname{Spec}R'.
 \end{aligned}
\]
The chosen pointed strict henselizations have residue $k$.
The line bundle $L'$ is ample. The normal full-section-ring and
proper-contraction argument of C.3as, Step~3, applies verbatim:
the contraction is an isomorphism off the zero section, using a
projectively normal Veronese and two consecutive sufficiently large
powers. Degree-one generation is not assumed. Thus $R'$ is excellent
normal of dimension $\dim V'+1=m+r+1=N$, $Z'$ is regular and proper
over $R'$, its reduced exceptional divisor is $i':V'\hookrightarrow Z'$,
and $X'=Z'\setminus V'$ is the regular puncture of $R'$.
Write $\pi':Z'\to V'$ for the flat projection.

\emph{Step 2: an actual flat morphism on the resolutions.}
Multiplication by the section $z_0$ defines a bundle morphism
\[
 \phi_0:Z'_0\longrightarrow Z_0
\]
over $V$; it sends a vector in
$\mathcal O(-1)\boxtimes L^{-1}$ to its product with $z_0$ in
$L^{-1}$. On the full section rings its contraction is the graded
map
\[
                  A_n\longrightarrow A'_n,\qquad
                               a_n\longmapsto z_0^n a_n.
\]
It sends the vertex to the vertex and gives a local homomorphism
$A_{A_+}\to R'$. Since $R'$ is strictly henselian with the chosen
residue field, the pointed universal property extends this to
$R\to R'$. The resulting compatible square defines
\[
                          \phi:Z'\longrightarrow Z.
\]

We verify flatness over arbitrary coefficient rings, not by a
torsion-free criterion valid only over a field. Trivialize $L$ on
an affine open of $V$, with coordinate ring $B$. On the chart
$z_0\ne0$, the map on fibre coordinates is the identity and the
remaining $m$ projective coordinates are polynomial variables.
On a chart $z_j\ne0$, $j\ne0$, write $u=z_0/z_j$ and let $t'$
be the corresponding fibre coordinate. Apart from $m-1$ extra
polynomial variables, the ring map is
\[
                  B[t]\longrightarrow B[u,t'],\qquad t\longmapsto ut'.
\]
The target is a free module over $B[t]$, with basis
\[
                    1,\ u^a\ (a\ge1),\ (t')^b\ (b\ge1).
\]
Indeed each monomial $u^a(t')^b$ is uniquely a power of $ut'$
times one of these basis monomials; their exponent pairs are
disjoint, so the expression is unique over any $B$. This proves
that $\phi_0$ is flat. Its fibres have pure dimension $m$, including
the fibre $ut'=0$.

The map $Z'\to Z'_0$ is ind-etale, so the composite
$Z'\to Z_0$ is flat. Factoring through the ind-etale map
$Z\to Z_0$ preserves flatness. Here is a direct justification.
For an etale $B$-algebra $C$ and a $B$-flat $C$-algebra $D$,
the graph of $\operatorname{Spec}D\to\operatorname{Spec}C$ is
an open immersion into
$\operatorname{Spec}D\times_B\operatorname{Spec}C$; projection
to $\operatorname{Spec}C$ is flat. Hence $D$ is $C$-flat.
For $C=\varinjlim C_\lambda$ ind-etale, this applies to every
etale stage. The finite-relations criterion for module flatness
then gives flatness over $C$, because the finitely many coefficients
in a relation descend to one stage. This argument applies on the
displayed schemes and proves that $\phi$ is flat. We do not assert
that an arbitrary ind-etale localization preserves all fibre
dimensions; the cycle pullbacks used below are defined at finite
stages and have their codimension checked by regular sequences.

There is deliberately no asserted morphism $X'\to X$:
over $z_0=0$, a nonzero fibre vector can map to the exceptional
divisor of $Z$. We pull cycles back on $Z$, and only subsequently
pass to the quotient by the exceptional image on $Z'$.

\emph{Step 3: the integral conormal detector in codimension $r$.}
The integral Milnor Bloch formula
\citep[Theorems~7.1 and 7.5]{cyKerz2009Milnor} applies to the
regular noetherian $Z'$ containing the infinite field $k$.
Apply its identification
$\mathrm{CH}^r(Z')=H^r(Z',\mathcal K_r^M)$, the symbol map
$d\log:\mathcal K_r^M\to\Omega^r_{Z'/k}$, coherent restriction
to $V'$, and projection using the split line-bundle conormal sequence.
This gives
\[
 \delta':\mathrm{CH}^r(Z')\longrightarrow
          H^r(V',L'\otimes\Omega^{r-1}_{V'/k}).
\]
The splitting is
\[
 i'^*\Omega^r_{Z'/k}
  =\Omega^r_{V'/k}\oplus
                  (L'\otimes\Omega^{r-1}_{V'/k}).
\]
As in C.3as, coherent restriction retains the conormal generator
$dt'$; it is not the differentials map that kills that generator.

Project the second summand further onto the direct summand with
all its $r-1$ horizontal differential factors from $V$.
Coherent K\"unneth and the vanishing of higher cohomology of
$\mathcal O_{\mathbb P^m}(1)$ give
\[
 H^r\bigl(V',\mathcal O(1)\boxtimes
                          (L\otimes\Omega^{r-1}_{V/k})\bigr)
   =H^0(\mathbb P^m,\mathcal O(1))\otimes
                            H^r(V,L\otimes\Omega^{r-1}_{V/k}).
\]
This is the ordinary coherent formula
\citep[Lemma 33.29.1, Tag 0BED]{D_StacksSegreCohomology}.
Pair the second factor with $\eta$ and apply the Serre trace on $V$.
On the first factor use the linear functional
$\epsilon(\sum a_jz_j)=a_0$. The resulting group homomorphism is
\[
                          \rho':\mathrm{CH}^r(Z')\longrightarrow k.
\]

It kills the full exceptional image. For every
$\beta\in\mathrm{CH}^{r-1}(V')$, the actual total-space
zero-section identity and flat base change to $R'$ give
\[
                  i'_*\beta=\pi'^*(-c_1(L')\cdot\beta)
                                           \quad\text{in }\mathrm{CH}^r(Z').
\]
The reduced zero-section square is Cartesian after the pointed
strict-local base change, just as in C.3as. The right side is
pulled back from $V'$ and hence has zero conormal $d\log$ component.
Thus localization descends $\rho'$ to $\mathrm{CH}^r(X')$.
The ordinary subgroup is divisible by $p$: compactify each actual
parameter curve, close the horizontal cycles, and use the divisible
group $\operatorname{Pic}^0$ of the smooth proper parameter curve,
exactly as in C.3as, Step~5. Since $k$ has exponent $p$ additively,
$\rho'$ kills that subgroup. It therefore factors through $B^r(X')$.

\emph{Step 4: actual cycles attain every scalar.}
For $b\in k$, take the flat pullback on the resolutions
\[
                    \xi_b=\phi^*([D_b]-[D_0])
                                           \in\mathrm{CH}^r(Z').
\]
Flat pullback here can be defined at finite pointed-etale stages
and then passed to the strict-local limit, as in C.3e; each original
curve and its finite equations descend to one such stage.
Equivalently, the selected curve $D_b$ is a regular curve: its
finite local normalization is a normal one-dimensional local
scheme, and its embedding is the isomorphism proved in C.3as,
Step~6. Its closed immersion into the regular $Z$ is a regular
immersion of codimension $r$. Flat pullback preserves that regular
immersion and its rank-$r$ conormal bundle. The resulting fundamental
cycles are therefore actual cycles of pure codimension $r$.
Their supports
meet $V'$ only along $\mathbb P^m\times\{v\}$, since
$\phi|_{V'}$ is projection to $V$ and the original selected curves
meet $V$ only at $v$.

We compute their supported cycle classes directly, so no additional
flat-functoriality claim for an unspecified cohomology theory is
needed. On the chart $z_j\ne0$, put $u_j=z_0/z_j$; for $j=0$
this is $1$. The pullback of the local regular sequence for $D_b$
is
\[
                  s_1,\ldots,s_{r-1},z-bu_jt'_j.
\]
It is again a regular sequence, with multiplicity one, because its
last differential has coefficient one in the etale coordinate $z$.
The supported integral Milnor class is the successive boundary
of its symbol. The product and residue identities in
\citet[Sections~3--4, especially equation~(24)]{cyBraunling2013}
give the corresponding supported Hodge class as the wedge of the
differentials divided by the product of the equations, as explicitly
proved in C.3as, Step~7.

Subtract the sequence for $b=0$ and restrict to $t'_j=0$, retaining
$dt'_j$. The term $t'_jdu_j$ vanishes, and the conormal part is
\[
 \pm b\,u_jdt'_j\otimes
       \left[\frac{ds_1\wedge\cdots\wedge ds_{r-1}}
                         {s_1\cdots s_{r-1}z}\right].
\]
The coefficient $u_j$ is exactly the local representative of the
global section $z_0$ of $\mathcal O(1)$. The local classes therefore
glue to $z_0$ tensored with the original supported class on $V$.
Equivalently, at the zero section the differential of the actual
bundle map is $dt=z_0dt'$, with no surviving projective differential
term. The product \v Cech description of coherent K\"unneth makes
the displayed identification compatible with local support and its
map to global cohomology.

Pairing with $\eta$, the Serre trace on $V$ gives $\pm b$, with
coefficient one, by the fixed normalization in C.3as. The functional
$\epsilon$ sends $z_0$ to $1$. Choose the same wedge orientation
as in that result. There are no other exceptional-intersection
contributions, so
\[
                              \rho'(\xi_b)=b.
\]
Every scalar is attained on an actual cycle. After restricting the
cycle to $X'$ and passing to the ordinary quotient, this proves the
asserted surjection.
\end{proof}

This extension uses only the explicitly proved C.3as projective
tower and actual curves. Its conclusion is integral and concerns
ordinary algebraic equivalence. It is compatible with rational
cone rigidity: the exhibited additive quotient disappears after
tensoring with $\mathbb Q$. It supplies no higher ordinary
restriction kernel or cokernel, because its detector kills the
ordinary subgroup itself.
